\section{Exact checks for odd standard polynomials}\label{app:verification}
The programs \path{code/verify_standard.py} and
\path{code/graded.py} provide finite exact checks of the odd
path, cycle, periodic and exceptional-order calculations.
The first uses SymPy and evaluates alternating products by a
full-matrix subset recurrence and by enumerating nonzero
matrix-unit Euler trails. The second computes cyclic-homogeneous
slice restrictions by a graded subset recurrence.

For fixed inputs $A_1,\ldots,A_m$, let $D(S)$ denote the alternating
sum over the variables indexed by $S\subseteq\{1,\ldots,m\}$.
The full-matrix recurrence is
\begin{equation}\label{eq:dp}
 D(\varnothing)=I,\qquad
 D(S)=\sum_{j\in S}(-1)^{|S|-\operatorname{pos}_S(j)}
 D(S\setminus\{j\})A_j,
\end{equation}
where $\operatorname{pos}_S(j)$ is the position of $j$ in the
increasing listing of $S$. For homogeneous cyclic inputs
$A_j=\diag(w_j)P^{g_j}$, the graded calculation tracks the
offset of each product and constructs the restriction matrix
directly. The two calculations compare every column
$L(E_{i,i+h\bmod n})$ for the specified exceptional cases;
the generated JSON also includes integer determinants,
finite trail counts, periodic pairs and metadata for sample Jordan
constructions. The examples cover orders $3$ through $8$, with
direct full-alternator evaluation checked for orders $3$ through $5$.

With Python 3.10 or later and the packages in
\path{requirements.txt}, run from the manuscript directory:
\begin{samepage}
\begin{verbatim}
python -m pip install -r requirements.txt
python code/verify_standard.py --output standard_rerun.json
python code/verify_standard.py --extended \
  --output standard_extended_rerun.json
\end{verbatim}
\end{samepage}
Compare the parsed outputs with
\path{certificates/standard_certificates.json} and
\path{certificates/standard_certificates_extended.json}, respectively.
The tests use assertions, so \texttt{python -O} must not be used.
The corresponding finite ranges are described in
\path{REPRODUCIBILITY.md}.

\section{Even-standard integer certificates}\label{app:even-verification}
The seven even-core certificates, indexed by their pairs $(r,N)$,
are the following files:
\begin{center}
\small
\begin{tabular}{@{}l@{}}
\path{certificates/border_r4_N5.json}\\
\path{certificates/border_r4_N6.json}\\
\path{certificates/border_r4_N7.json}\\
\path{certificates/border_r5_N6.json}\\
\path{certificates/border_r5_N7.json}\\
\path{certificates/border_r5_N8.json}\\
\path{certificates/border_r5_N9.json}
\end{tabular}
\end{center}
Each file contains integer matrices $B_1,\ldots,B_{2r-1}$,
a column $u$, the nonzero integer determinant of
$S=s_{2r-1}(B)$, and a specified $N^2\times N^2$ minor of
the normal map $\Gam_{B,u}$ with its integer determinant.
The nonzero determinant of $S$ and minor of $\Gam$ establish
admissibility \eqref{eq:core-admissible} for that core size.

The constructor \path{code/verify_border_even.py} uses an
alternating-product subset recurrence and replaced-input
columns, with direct word enumeration for smaller cases.
The saved-input checker \path{code/verify_saved_border.py}
uses a different first-letter recurrence to rebuild every
normal-map column; it checks $\det S$, the bordered rank-one
square-zero value and the recorded minor modulo $101$ and
$103$. A nonzero residue of the integer minor establishes
its nonvanishing over characteristic-zero fields.
The seed checker \path{code/verify_minimal_even.py} tests the
signed matrix-unit formulas and full differential in specified
finite sizes.

For example, run:
\begin{verbatim}
python code/verify_saved_border.py \
  --input certificates/border_r4_N5.json \
  --output checked_border_r4_N5.json
python code/verify_minimal_even.py \
  --max-n 8 --rank-max-n 6 \
  --output minimal_even_rerun.json
\end{verbatim}
The first command applies likewise to the other six certificate
files; \path{REPRODUCIBILITY.md} gives the full list and
constructor command. These integer checks establish the listed
finite admissible cores. Their extension to infinitely many
matrix orders is the symbolic additive transfer
of Theorem~\ref{thm:core-add}, while the general even-degree
argument uses the cited generic-nullity theorems.
